% BHU PYQ -- Manually transcribed & error-corrected
% Paper: GRA-101: Physical Basis of Geography
% Exam: B.A. (Hons.) Semester I Examination, 2022-23

\documentclass[11pt,a4paper]{article}
\usepackage[a4paper,top=2.5cm,bottom=2.5cm,left=2.8cm,right=2.8cm,headheight=14pt]{geometry}
\usepackage{amsmath,amssymb}
\usepackage[T1]{fontenc}
\usepackage[utf8]{inputenc}
\usepackage{lmodern,microtype,fancyhdr,enumitem,hyperref,lastpage,parskip}

\hypersetup{colorlinks=true,urlcolor=black,linkcolor=black,
  pdftitle={GRA-101 Physical Basis of Geography -- BHU B.A. Sem I 2022-23}}
\pagestyle{fancy}\fancyhf{}
\renewcommand{\headrulewidth}{0.4pt}\renewcommand{\footrulewidth}{0.4pt}
\fancyhead[L]{\small   \textit{Banaras Hindu University}}
\fancyhead[R]{\small   \textit{GRA-101: Physical Basis of Geography}}
\fancyfoot[C]{\small Page  \thepage\ of \pageref{LastPage}}


\newlist{parts}{enumerate}{2}
\setlist[parts,1]{label=(\alph*),leftmargin=*,itemsep=5pt,topsep=3pt}

\newcommand{\pts}[1]{\hfill   \textit{[#1]}}


\begin{document}

\begin{center}
  {\large\bfseries BANARAS HINDU UNIVERSITY}\\[2pt]
  {\normalsize B.A. (Hons.) --- Semester I Examination, 2022-23}\\[6pt]
  \rule{0.8\linewidth}{0.5pt}\\[6pt]
  {\large\bfseries Paper No. GRA-101: Physical Basis of Geography}\\[4pt]
  {\normalsize Time: 3 Hours \hfill Maximum Marks: 70}\\[2pt]
  {\small Ref. No.: 054440}
\end{center}
\rule{\linewidth}{0.4pt}
\smallskip



\noindent   \textit{Instructions: Answer   \textbf{five} questions in all. Question No. 1, carrying 20 marks, is compulsory. Remaining questions are of equal marks.}
\bigskip




\noindent   \textbf{Question 1.} Write short answer of the following: \pts{20}

\begin{parts}
  \item Describe briefly the solar system. \pts{4}
  \item Discuss in brief volcanoes with examples. \pts{4}
  \item Describe Indian Ocean currents. \pts{4}
  \item Cyclones and Anticyclones. \pts{4}
  \item Explain briefly the relevance of rotation. \pts{4}
\end{parts}

\medskip


\noindent   \textbf{Question 2.} Describe the composition and structure of the atmosphere. \pts{12.5}

\medskip


\noindent   \textbf{Question 3.} What is weathering? Give a detailed account of various types of weathering. \pts{12.5}

\medskip


\noindent   \textbf{Question 4.} Present a detailed account of the landforms of river erosion and deposition with a well-labeled diagram. \pts{12.5}

\medskip


\noindent   \textbf{Question 5.} What is insolation? Discuss the factors influencing the distribution of insolation. \pts{12.5}

\medskip


\noindent   \textbf{Question 6.} Discuss the tetrahedral theory of the origin of the ocean basin. \pts{12.5}

\medskip


\noindent   \textbf{Question 7.} Explain the distribution of salinity and temperature in oceanic water. \pts{12.5}

\medskip


\noindent   \textbf{Question 8.} Explain the binary star hypothesis of Russell associated with the origin of the earth. \pts{12.5}

\vfill\rule{\linewidth}{0.4pt}

\begin{center}\small
\href{https://bhu-exam.vercel.app/}{Click here to visit our website}\\[4pt]
\href{https://me-aryan.vercel.app/}{Click here to visit our contributor}\\[4pt]
\href{https://quashan-me.vercel.app/}{Click here to visit our contributor}
\end{center}
\end{document}
